\section{Descripción del caso de uso}

\begin{resumencaso}
\textbf{Propósito del modelo.} El modelo Entidad--Relación Extendido de
\textit{PuellaGame} representa la información necesaria para seguir la
experiencia de un cliente dentro de un centro de entretenimiento familiar y,
al mismo tiempo, sostener la operación cotidiana de sus sucursales. El alcance
abarca desde la adquisición y uso de una tarjeta hasta las partidas, la
acumulación de puntos, los canjes, el inventario de premios, el mantenimiento
de las máquinas y la participación del personal.
\end{resumencaso}

\subsection{Contexto del sistema}

PuellaGame es una organización de entretenimiento presencial distribuida en
sucursales. Cada sucursal reúne máquinas de juego, un catálogo físico de
premios y personal con responsabilidades diferenciadas. La operación no se
reduce al instante en que una persona juega: comienza cuando se registra a un
cliente y se le asocia una tarjeta, continúa con las recargas y las partidas,
y se extiende hasta la conversión de los puntos obtenidos en premios. Por ello,
el modelo enlaza la actividad recreativa con los procesos administrativos que
permiten que el servicio funcione.

El proyecto surge porque la información de las distintas sucursales se mantiene
en hojas de cálculo y registros físicos. El crecimiento de la organización hace
necesario sustituir esa administración fragmentada por una base de datos que
concentre los registros, preserve su consistencia y permita consultar la
operación de las sucursales de manera uniforme.

La tarjeta constituye el medio operativo del cliente. En ella se concentra el
saldo disponible para jugar y el total de puntos acumulados; también permite
relacionar cada partida con la persona que la realizó. Una partida ocurre en
una máquina determinada y conserva su fecha, hora, puntaje y puntos ganados.
Las máquinas corresponden a un juego y pertenecen a una sucursal, de manera
que el sistema puede reconocer qué experiencia se ofreció, en qué equipo y en
qué ubicación.

El modelo distingue cuatro modalidades de juego: destreza, disparos, velocidad
y ritmo. Esta especialización permite expresar reglas de recompensa propias de
cada modalidad sin perder los datos comunes de todo juego. La bonificación
asociada a la categoría VIP del cliente se conserva como un valor derivado en
la partida, pues depende tanto de la actividad registrada como de la condición
vigente del cliente.

\Needspace{10\baselineskip}
\subsection{Flujos cubiertos}

El primer flujo corresponde al acceso y consumo del servicio. El sistema
registra la visita de un cliente a una sucursal, la recarga recibida por su
tarjeta y las partidas ejecutadas en las máquinas disponibles. Con esta
secuencia se conserva una trazabilidad básica entre cliente, tarjeta, sucursal,
juego y resultado, sin convertir el modelo conceptual en una descripción de
pantallas o de procedimientos de software.

Para ingresar a este flujo, el jugador compra una tarjeta por \$20.00 MXN y
realiza una recarga inicial de \$10.00 MXN. La tarjeta permite pagar las
partidas y conservar digitalmente los puntos ganados. Cada jugador puede tener
una sola tarjeta activa; las recargas posteriores se registran junto con su
monto, fecha, hora, sucursal y medio de abono.

El segundo flujo corresponde a la fidelización. Los puntos obtenidos mediante
las partidas se acumulan en la tarjeta y posteriormente pueden utilizarse en
un canje. Cada canje se vincula con el cliente que lo realiza, con la sucursal
en la que ocurre y con el premio otorgado. La disponibilidad de un premio se
entiende por sucursal, ya que un mismo artículo puede existir en cantidades
distintas según la ubicación. Esta separación permite representar el inventario
local sin confundirlo con las propiedades generales del premio.

El tercer flujo describe la continuidad operativa. Las máquinas reciben
mantenimientos que son realizados por personal técnico; a su vez, el personal
se adscribe a una sucursal y se especializa, de manera disjunta, en gerente,
cajero, encargado de premios o técnico. Los cajeros intervienen en las
recargas, los encargados de premios procesan los canjes, los técnicos atienden
el mantenimiento y los gerentes administran la sucursal. De este modo, el
modelo conserva la responsabilidad humana de las operaciones relevantes.

\subsection{Reglas del negocio}

\begin{resumencaso}
Las reglas siguientes delimitan el significado de los registros del modelo.
No describen pantallas ni algoritmos: expresan condiciones del negocio que una
implementación posterior deberá preservar.
\end{resumencaso}

\begin{description}
  \item[Tarjetas, costos y recargas.]
  Cada jugador mantiene como máximo una tarjeta activa. Tanto el costo de una
  partida como el importe de una recarga son múltiplos de diez. El método de
  abono de una recarga puede ser efectivo, tarjeta de crédito o débito, o
  PayPal.

  \item[Visitas y categoría VIP.]
  Una visita se reconoce cuando el jugador usa su tarjeta en una sucursal
  durante una fecha determinada; varias partidas en esa misma sucursal y día
  constituyen una sola visita. Al acumular al menos diez visitas en un mes, el
  jugador obtiene la categoría VIP durante el mes siguiente y recibe un
  diez por ciento adicional sobre los puntos de cada partida. Para conservar
  el beneficio debe cumplir nuevamente el mínimo mensual.

  \item[Recompensas por tipo de juego.]
  Los juegos de destreza y disparos conceden un punto por cada caso de éxito.
  Los juegos de velocidad y ritmo otorgan veinte puntos cuando el resultado
  queda entre los primeros cinco lugares. La partida conserva, además, el
  puntaje obtenido, los puntos ganados y el momento en que terminó.

  \item[Premios y canjes.]
  Los premios se orientan a los rangos infantil (de 3 a 12 años), juvenil (de
  3 a 17 años) y adulto (18 años o más). Por su costo en puntos se clasifican
  como bajos, de 20 a 1,000 puntos; medios, de 1,001 a 3,999; y grandes, de
  4,000 a más de 10,000 puntos. Un canje requiere al menos veinte puntos y el
  valor aproximado del premio equivale al cincuenta por ciento de los puntos
  utilizados expresados en pesos; por ejemplo, veinte puntos permiten obtener
  un premio cercano a \$10.00 MXN.

  \item[Inventario y responsabilidad operativa.]
  La cantidad disponible de cada premio se controla por sucursal. El cajero
  vende tarjetas y refleja recargas; el encargado de premios verifica puntos,
  procesa el canje y entrega el premio; el técnico registra los mantenimientos
  de las máquinas; y el gerente administra la sucursal.
\end{description}

\subsection{Alcance y frontera conceptual}

\begin{center}
\begin{tabularx}{0.94\textwidth}{@{}>{\bfseries\color{unamblue}}p{0.25\textwidth}X@{}}
\toprule
Dentro del alcance & Clientes y tarjetas; visitas; recargas; partidas y
puntajes; juegos y máquinas; sucursales; premios, existencias y canjes;
mantenimientos; empleados y funciones operativas. \\
\midrule
Fuera del alcance & Contabilidad, proveedores, compras, nómina, reservaciones,
mercadotecnia, autenticación, infraestructura tecnológica y detalles de
implementación de la aplicación. \\
\bottomrule
\end{tabularx}
\end{center}

La frontera anterior mantiene el diagrama centrado en la información estable
del negocio. Los elementos incluidos son aquellos que permiten reconstruir el
ciclo principal de PuellaGame y asignar responsabilidad a sus operaciones. Los
procesos excluidos podrían incorporarse en versiones posteriores, pero no son
necesarios para explicar el dominio representado en esta primera versión del
modelo Entidad--Relación Extendido.
